% 原書379--383頁、Pauli--Villars regulator から super-Riemann surface まで。
\gloss{Pauli--Villars正則化（Pauli--Villars regulator）}
{通常とは逆の統計に従う、非常に重い場を導入して量子場理論を正則化する方法。}

\gloss{Peccei--Quinn対称性、PQ対称性（Peccei--Quinn symmetry）}
{アノマリーによってのみ破られる近似的対称性。}

\gloss{周期行列（period matrix）}
{種数$g$のRiemann面の複素構造を特徴付ける$g\times g$行列。}

\gloss{摂動論（perturbation theory）}
{量子系の振幅を、結合のべきで展開すること。}

\gloss{物理状態（physical state）}
{拘束条件を持つ量子系において、拘束条件によって消される状態。OCQでは、Virasoro下降生成子によって消され、指定された$L_0$固有値を持つ状態である。BRST量子化では、BRST演算子によって消される状態である。どちらの場合も、真の物理スペクトルは物理状態の空間にさらに同値関係を課したもので、ヌル状態だけ異なる物理状態を同一視する。}

\gloss{ピクチャー（picture）}
{RNS超弦における、頂点演算子のいくつかの同型な表現の1つ。$q$ピクチャーは、$\beta\gamma$ゴースト電荷が$q$の頂点演算子からなる。自然なピクチャーは$q=-1$および$-\frac12$であり、より高いピクチャーには超モジュライ空間についての部分的な積分が含まれる。ピクチャー変更演算子は$q$を1増やす。}

\gloss{Planck長（Planck length）}
{$\hbar$、$c$、$G_{\mathrm N}$から構成される、量子重力の自然な長さスケール。
$L_{\mathrm P}=M_{\mathrm P}^{-1}=1.6\times10^{-33}\,\mathrm{cm}$。}

\gloss{Planck質量（Planck mass）}
{$\hbar$、$c$、$G_{\mathrm N}$から構成される、量子重力の自然な質量スケール。
$M_{\mathrm P}=1.22\times10^{19}\,\mathrm{GeV}$。}

\gloss{配管接合（plumbing fixture）}
{取っ手を縫い付けることによって、低い種数のRiemann面から高い種数のRiemann面を構成する手続き。この構成はパラメータ$q$を含み、$q$が0になると取っ手が退化、すなわちくびれる。Riemann面のモジュライ空間の境界の標準的な表示を与える。}

\gloss{Poincaré双対（Poincaré dual）}
{完全反対称テンソルとの縮約によって与えられる、$p$形式から$(d-p)$形式への写像。}

\gloss{Poincaré不変性（Poincaré invariance）}
{平坦な計量の不変群であり、並進とLorentz変換からなる。}

\gloss{点群（point group）}
{並進を無視したオービフォールド群$H$。時空的解釈を持つオービフォールドにのみ適用される。}

\gloss{Polyakov経路積分（Polyakov path integral）}
{第一量子化された弦理論を、独立な世界面計量を持つ経路積分として表す方法。局所Weyl対称性により、古典的自由度はNambu--Goto理論と同じであることが保証される。}

\gloss{一次場（primary field）}
{CFTにおいて、Virasoro代数などの与えられた代数のすべての下降生成子によって消される局所演算子。対応する状態は最高ウェイト状態である。}

\gloss{射影空間（projective space）}
{線形な$(n+1)$次元空間の点を、一様なスケール変換
$(x_1,\ldots,x_{n+1})\cong(\lambda x_1,\ldots,\lambda x_{n+1})$
の下で同一視して構成される、コンパクトな$n$次元空間。
$x_i$と$\lambda$が実数なら$\mathbf{RP}^n$、複素数なら$\mathbf{CP}^n$を生む。後者は$n$複素次元、すなわち$2n$実次元を持つ。}

\gloss{擬スピン（pseudospin）}
{カレント代数において、重心モードのLie代数には含まれない$SU(2)$部分代数。}

\gloss{穿孔（puncture）}
{Riemann面上の印を付けた点であり、頂点演算子の位置。}

\gloss{QCD弦（QCD string）}
{非可換ゲージ理論の弦理論としての再定式化。少なくとも色の数が大きい極限では存在すると予想されている。}

\gloss{四元数多様体（quaternionic manifold）}
{ホロノミー群$Sp(k)\times SU(2)\subset SO(4k)$を持ち、$SU(2)$曲率が特定の形を取る$4k$次元多様体。$d=6,\ N=1$または$d=4,\ N=2$超重力のハイパー多重項のモジュライ空間の幾何である。}

\gloss{$R$対称性（R symmetry）}
{超カレントに非自明に作用する対称性。}

\gloss{上昇演算子（raising operator）}
{与えられた状態のエネルギーを下げる演算子。CFTでは、Virasoro生成子$L_0$または$\widetilde L_0$を$n$だけ増やす演算子は、次数、すなわち下付き添字$-n$を持つ。}

\gloss{Ramond代数（Ramond algebra）}
{超カレントが周期的であり、そのモードが整数値となるセクターで、超カレントとエネルギー運動量テンソルのFourierモードがなす世界面の代数。}

\gloss{Ramond境界条件、R境界条件（Ramond boundary condition）}
{閉弦、または開弦を二重化したものの世界面で、フェルミオン場が周期的であるという条件。二重化の方法を参照。このときFourierモードは整数値となる。}

\gloss{Ramond--Neveu--Schwarz超弦、RNS超弦（Ramond--Neveu--Schwarz superstring）}
{超共形不変性を持つが時空超対称性は顕在的でない、I型およびII型超弦の定式化。時空超対称性は、弦のHilbert空間にGSO射影を課すと現れる。}

\gloss{Ramond--Ramond状態、R--R状態（Ramond--Ramond states）}
{I型およびII型超弦理論における、左向き・右向きの部分がともにフェルミオン的である、ボソン的な閉弦状態。$p$形式ポテンシャル$C_p$を含み、IIA弦では$p$はすべての奇数値、IIB弦ではすべての偶数値を取る。}

\gloss{階数（rank）}
{1．Lie代数の、互いに可換な生成子の最大個数。2．テンソルの添字の数。}

\gloss{有理CFT（rational CFT）}
{拡張されたカイラル代数の下で有限個の一次場を持つCFT。極小模型の一般化である。このようなCFTには強い制約がある。}

\gloss{再フェルミオン化（refermionization）}
{ボソン化の後で、ボソン場の線形結合から新たなスピン$\frac12$の場を構成すること。これらのフェルミオンは、元のフェルミオンの非局所的・非線形な関数である。}

\gloss{Regge挙動（Regge behavior）}
{重心エネルギーの2乗$s$が大きく、運動量移行の2乗$-t$を固定した散乱において、振幅が$s^{\alpha(t)}$としてスケールすること。$\alpha(t)=j$が非負整数となる$t$の値は、スピン$j$、質量の2乗$-t$の粒子の交換に対応する。}

\gloss{Regge勾配（Regge slope）}
{$\alpha'$で表される、摂動的弦理論の特徴的な長さスケールの2乗。基本弦の張力は$1/(2\pi\alpha')$である。}

\gloss{関連相互作用（relevant interaction）}
{距離とともに無次元の強さが増大する相互作用。摂動論では、超繰り込み可能な相互作用に等しい。}

\gloss{繰り込み群方程式（renormalization group equation）}
{長さの尺度による物理の変化を支配する微分方程式。}

\gloss{繰り込み理論（renormalization theory）}
{経路積分の微積分学。}

\gloss{Riemann--Rochの定理（Riemann--Roch theorem）}
{Riemann面上の計量モジュライの数から共形Killingベクトルの数を引くと$-3\chi$になるという定理。$\chi$は曲面のEuler数である。また、この結果の一般化も指す。}

\gloss{Riemann面（Riemann surface）}
{2実次元の複素多様体。Riemann多様体のWeyl同値類と等価である。}

\gloss{ルート（root）}
{随伴表現の状態に対応する、Lie代数の互いに可換な生成子の最大集合の固有値を並べたベクトル。}

\gloss{$S$双対性（S-duality）}
{量子論の結合定数が非自明に変わる双対性。弱結合・強結合双対性を含む。重要な例は、IIB弦理論、および$d=4,\ N=4$超対称Yang--Mills理論の$SL(2,\mathbf Z)$自己双対性である。より広い意味では、IIA理論と円周上のM理論の双対性、$SO(32)$ヘテロティック理論とI型理論の双対性、$E_8$ヘテロティック理論と区間上のM理論の双対性など、異なる理論間の弱結合・強結合双対性にも使われる。コンパクト化した理論では、$S$双対性という用語は、結合に依存する一様なスケール変換を除いて半径を不変に保つ双対性に限られる。$T$双対性、$U$双対性と対比されたい。}

\gloss{$S$行列（S-matrix）}
{無限の過去の状態と無限の未来の状態の重なりの振幅、すなわち散乱振幅。座標不変な量子論では、一般に最も単純な不変量である。この用語は通常、自由粒子状態の基底を前提とする。質量のない粒子を含む理論ではIR発散のために問題が生じ、非自明なIR固定点にある理論では意味を持たない。}

\gloss{スケール変換（scale transformation）}
{時空または世界面の一様なスケールの変更。}

\gloss{Scherk--Schwarz機構（Scherk--Schwarz mechanism）}
{時空の回転を含む次元縮約による超対称性の破れ。}

\gloss{Schottky群（Schottky group）}
{球面$S_2$の$SL(2,\mathbf C)$ Möbius変換の離散部分群$\Gamma$で、$S_2/\Gamma$が多様体、特にRiemann面となるような付加条件を満たすもの。}

\gloss{Schwarz微分（Schwarzian）}
{エネルギー運動量テンソルの有限共形変換に現れる、微分の組合せ。}

\gloss{Schwinger--Dyson方程式（Schwinger--Dyson equation）}
{量子論の演算子方程式を、期待値の方程式として表したもの。}

\gloss{短い多重項（short multiplet）}
{BPS状態を参照。}

\gloss{シグマ模型計量（sigma model metric）}
{弦の世界面作用に現れる計量。弦計量とも呼ばれ、Einstein計量とはディラトンに依存するWeyl変換だけ異なる。}

\gloss{単純カレント（simple current）}
{任意の一次場$\mathcal O$に対して、演算子積$J\mathcal O$がただ1つの共形族だけを含むような演算子$J$。}

\gloss{ソリトン（soliton）}
{古典極限が、滑らかで局在した、通常は位相的に非自明な古典場の配位となる状態。すべての方向に局在する粒子状態だけでなく、広がった対象も含む。一方、通常の量子の状態は、古典極限の近くで定数配位のまわりの小さな揺らぎとして表される。複数の古典極限、すなわち双対性を持つ理論では、ソリトンと量子の役割が入れ替わることがある。}

\gloss{特殊Kähler幾何（special Kähler geometry）}
{$d=4,\ N=2$超重力のベクトル多重項のモジュライ空間の幾何。シンプレクティック変換の任意性を除いて定まる特殊座標によって最も簡単に定義される（第B.7節）。重力が切り離され、超対称性が大域的になる極限では、剛体特殊幾何が得られる。}

\gloss{スペクトル流（spectral flow）}
{境界条件の連続的な変化によって生じる、スペクトルの断熱的変化。}

\gloss{スピン（spin）}
{回転の下での場または状態の振る舞い。CFTでは、共形ウェイトによって$h-\widetilde h$と与えられる。}

\gloss{スピン場（spin field）}
{Ramond基底状態の頂点演算子。スピノル場に分岐切断を生じさせる。}

\gloss{スピン構造（spin structure）}
{多様体上にスピノル場を大域的に定義する、互いに非同等ないくつかの方法の1つ。大まかには、遷移関数の平方根の符号の選択に対応する。}

\gloss{スプリアス状態（spurious state）}
{OCQにおいて、Virasoro上昇演算子によって作られる状態。}

\gloss{状態・演算子同型（state--operator isomorphism）}
{CFTにおいて、円周上で量子化した理論の状態と局所演算子との1対1の同型。また、区間上で量子化した理論の状態と境界上の局所演算子との1対1の同型。$d$次元では、円周は$(d-1)$次元球面、区間は$(d-1)$次元半球面となる。}

\gloss{弦結合（string coupling）}
{弦の摂動論で、異なるRiemann面の重みを決める無次元パラメータ$g$。Euler数$\chi$の曲面の寄与には$g^{-\chi}$の重みが付く。弦結合とディラトンの関係は$g=e^\Phi$である。この定義は閉弦を放出する振幅に対応し、開弦を放出する振幅は$g^{1/2}$に比例する。}

\gloss{弦の場の理論（string field theory）}
{弦理論を場の理論として表す方法。この場は、円周または区間から時空への写像であり、通常の量子場の無限個の組に対応する。この形式は弦の摂動論を再現できるが、摂動論を超えて定義できるかどうかは明らかでない。}

\gloss{弦計量（string metric）}
{シグマ模型計量を参照。}

\gloss{弦スケール（string scale）}
{弦の励起の系列を特徴付ける質量スケール$\alpha'^{-1/2}$。}

\gloss{弦・弦双対性（string--string duality）}
{異なる弦理論間の弱結合・強結合双対性、特に$T^4$上にコンパクト化したヘテロティック弦と、K3上にコンパクト化したIIA弦の双対性を表すために使われることのある用語。}

\gloss{弦の張力（string tension）}
{静止した弦の単位長さ当たりの質量。Regge勾配と$1/(2\pi\alpha')$の関係で結び付く。}

\gloss{Sugawara構成（Sugawara construction）}
{カレント代数において、2つのカレントの積としてエネルギー運動量テンソルを構成すること。もともとは4次元の現象論的模型として提案されたが、後に2次元では厳密な結果であることが分かった。}

\gloss{超共形代数（superconformal algebra）}
{反交換するスピノル生成子を含むように、共形代数、すなわちVirasoro代数を拡張したもの。$(N,\widetilde N)$超共形代数は、$N$個の左向き、$\widetilde N$個の右向きの超カレントを持つ。}

\gloss{超共形カレント代数（superconformal current algebra）}
{スピン$\frac32$とスピン1の両方のカレントを含むように共形変換を拡張したもの。}

\gloss{超共形場理論、SCFT（superconformal field theory）}
{超共形変換の下で不変な量子場理論。}

\gloss{超カレント（supercurrent）}
{保存するスピノルカレント。超共形変換に伴う世界面カレント$T_{\mathrm F}$と、時空超対称性に伴う時空カレントを含む。}

\gloss{超場（superfield）}
{座標変換の下で特定の変換性を持つ、超空間上の場。}

\gloss{超重力（supergravity）}
{一般相対論と超対称性の結合。超対称性を局所対称性へと昇格させることも意味する。}

\gloss{超多様体、超空間（supermanifold, superspace）}
{可換な座標と反可換なGrassmann座標の両方を含むよう、多様体の概念を形式的に拡張したもの。}

\gloss{超モジュライ（supermoduli）}
{超Riemann面を特徴付ける、反可換なパラメータ。}

\gloss{超対称パートナーのスケール（superpartner scale）}
{標準模型粒子の超対称パートナーの質量スケール。超対称性がHiggsスカラー質量の自然さの問題を解決するなら、$10^2$から$10^3\,\mathrm{GeV}$の間と期待される。}

\gloss{超ポテンシャル（superpotential）}
{$d=4,\ N=1$超対称性において、ゲージ相互作用以外の相互作用を決定する、超場の正則関数。}

\gloss{超Riemann面（super-Riemann surface）}
{座標近傍間の超共形な遷移関数によって定義される超多様体。}
